\section{Experimental setup}
\label{sec:experiments}

\subsection{Datasets and protocols}
\label{sec:experiments-datasets}

The primary benchmark is CIFAR-10 non-IID stationary (ResNet-18, 100 clients,
Dirichlet $\alpha{=}0.3$, no task structure, 200 rounds, three seeds). A
class-incremental stress test on the same scale (5 sequential 2-class tasks, no
replay) carries the 11-baseline positioning comparison, since unlearning is
hardest under catastrophic forgetting. Cross-dataset validation is on
Split-CIFAR-100 (non-IID stationary at the same hyperparameters) and FEMNIST
(natively non-IID by per-writer partition). Table~\ref{tab:datasets} lists all
benchmarks.

\paragraph{Extrapolation from non-privacy-relevant datasets to real-world privacy
settings.}
Split-CIFAR-10 and Split-CIFAR-100 are not privacy-sensitive datasets. We use them
because the privacy claim this paper makes is about the FL protocol and the
unlearning mechanism, not about the data distribution: the formal
$(\varepsilon, \delta)$-indistinguishability guarantee of
Theorem~\ref{thm:composition} is a property of the training trajectory and the head
subtraction, independent of what the images depict. The membership-inference
attacker (Yeom, LiRA) and the gradient-inversion attacker (FUIA) likewise operate
on model outputs and gradients, not on raw pixels; their advantages extrapolate to
sensitive data at the same model scale. CIFAR-scale image classification matches
the practical cross-device FL scenario of on-device vision models (keyboard
suggestion image attachments, smart-home object recognition, medical image triage)
in architecture (ResNet-18) and sampling regime (100 clients, Dirichlet
non-IID). FEMNIST is a privacy-relevant dataset by construction (per-writer
handwriting is identifiable); its qualitative agreement with the Split-CIFAR
results establishes that the protocol's behaviour is not an artefact of the
non-privacy benchmark.

\begin{table}[t]
\centering
\scriptsize
\setlength{\tabcolsep}{3pt}
\begin{tabular}{llrrrl}
\toprule
Dataset & Role & Cls & Tasks & Clients & Partition \\
\midrule
CIFAR-10            & primary       & 10  & 1  & 100   & Dir$(0.3)$ \\
Split-CIFAR-10      & stress test   & 10  & 5  & 100   & Dir$(0.5)$ \\
CIFAR-100           & cross-dataset & 100 & 1  & 100   & Dir$(0.3)$ \\
Split-CIFAR-100     & cross stress  & 100 & 10 & 100   & Dir$(0.5)$ \\
FEMNIST             & cross-dataset & 62  & 6  & 3{,}550 & writer \\
Split-Tiny-ImageNet & released code & 200 & 20 & 100   & Dir$(0.5)$ \\
20 Newsgroups       & released code & 20  & 4  & 100   & Dir$(0.3)$ \\
\bottomrule
\end{tabular}
\caption{Benchmarks. Primary and cross-dataset rows use non-IID stationary
partitioning with a single task; stress-test rows use the no-replay
class-incremental protocol.}
\label{tab:datasets}
\end{table}

\subsection{Training protocol}
\label{sec:experiments-protocol}

FORGE trains with Algorithm~\ref{alg:decayed-aggregation}; the retrain gold and the
shared base for baselines train with vanilla FedAvg. Table~\ref{tab:hyperparameters}
lists the CIFAR-10 hyperparameters, shared across the primary and stress-test
protocols (only the task count differs: 1 vs 5). The clip norm $C{=}50$ was
chosen after measuring local-update L2 norms in $[6, 60]$ for ResNet-18 at three
local epochs under Dirichlet partitioning; $\sigma{=}10^{-3}$ keeps the
Gaussian-mechanism magnitude comparable to the clipped update norm. CIFAR-100 uses
$C{=}1$, $\mathrm{lr}{=}0.01$, and 500 rounds to match the 100-class output
dimension; FEMNIST uses a small-CNN at $0.01$ learning rate for 500 LEAF writers.

\begin{table}[t]
\centering
\scriptsize
\setlength{\tabcolsep}{3pt}
\begin{tabular}{llr}
\toprule
Group & Hyperparameter & Value \\
\midrule
Federated & rounds $R$ / per-task / sample frac. & $200/40/0.20$ \\
          & local epochs / batch / LR & $3/32/0.02$ \\
          & clip norm $C$ & $50.0$ \\
\midrule
Decay   & $\alpha$ / server LR / server $\sigma$ & $0.5 / 1.0 / 10^{-3}$ \\
\midrule
Baselines & aggregator / server noise & FedAvg / $0$ \\
\midrule
Summary & type / freq. & moments / every 20 rounds \\
        & $\eps_\summary$ / $\delta_\summary$ & $2.0 / 10^{-5}$ \\
\midrule
Erasure & subspace / $\gamma$ / $R$ & head / $0.1 / 0.5$ \\
\midrule
Model   & architecture / output dim & ResNet-18 / $10$ \\
\bottomrule
\end{tabular}
\caption{CIFAR-10 hyperparameters.}
\label{tab:hyperparameters}
\end{table}

Departure is simulated at $r_\forget = 0$ (client~3 participates until the erasure
request at round 200), i.e., the maximum-influence case.

\subsection{Baselines}
\label{sec:experiments-baselines}

On the primary non-IID stationary benchmark, FORGE is compared against the
retrain gold, a no-unlearn control, and the two preserved-utility baselines
that reach retrain accuracy on the stress test (SIFU and
FedRecovery$_{\eps=\infty}$); the remaining three baselines (Halimi,
FedMemPrune, FedEraser) are omitted here because they collapse to chance on
the stress test and would be strictly dominated at realistic utility too.
On the class-incremental stress test, FORGE is compared additionally against
five FU methods: Halimi et al.\
\cite{halimi2022}, FedMemPrune \cite{wei2026fedmemprune}, FedRecovery
\cite{zhang2023fedrecovery} (swept over $\varepsilon \in \{1, 2, 5, 10, \infty\}$),
FedEraser \cite{liu2021federaser}, and SIFU \cite{fraboni2024sifu} with proper
fine-tuning replay. Baseline hyperparameters and provenance are in
Appendix~\ref{app:baseline-attribution}. For FORGE at $T=200$,
Theorem~\ref{thm:forward-forgetting} gives $\mu_1(T) \approx 0$ and the overall
guarantee is $\mu_\summary = 0.413$ (equivalently $\varepsilon_\summary = 2.0$).

\subsection{Metrics}
\label{sec:experiments-metrics}

\textbf{Utility.} Average test accuracy across all task-restricted test sets.
The gold retrain's accuracy sets the ceiling at the same training budget;
preserved-utility methods are those whose CIs overlap the retrain's.
\textbf{Forgetting quality.} The Yeom threshold MIA \cite{yeom2018privacy}
(advantage $|2\cdot\mathrm{bal.acc.} - 1|$, with departing-client training data as
member set); a shared-shadow LiRA \cite{carlini2022lira}; the retrain-KL
$\mathrm{KL}(\Prob_{\model_\mathrm{unlearn}} \| \Prob_{\model_\mathrm{retrain}})$ on a
256-sample probe; and the post-erasure accuracy, which disambiguates
destruction-induced ``forgetting''. FUIA reconstruction \cite{fuia2025} is reported
by PSNR.
\textbf{Cost.} Wall-clock latency of the erasure step and server-side storage
($O(1)$, $O(R|\theta|)$, or $O(NR|\theta|)$).

\subsection{Hardware and reproducibility}

All runs use an Apple M4 (MPS backend) with PyTorch 2.8. The primary non-IID
stationary benchmark (3 methods: retrain, no-unlearn, FORGE) completes in
approximately $20$~h per seed including LiRA at $n_\mathrm{shadow}{=}32$; the
3-seed total is approximately $60$~h. The 12-method class-incremental stress
test runs in approximately $5$~h per seed. Seeds are fixed, the training
loop snapshots every round, and the released code reproduces every benchmark
from a single command with deterministic resume across interruptions.
